\section{Setup: one coupling, one expansion}
\label{sec:2}

\subsection{The theory has one coupling, and it is a boundary observable}
\label{sec:2.1}

We consider Einstein gravity with a negative cosmological constant coupled to
$SU(2)$ Yang--Mills theory in five dimensions,
\begin{equation}
S = \frac{1}{2\kappa^2}\int d^5x\sqrt{-g}\,\Big[R + 12 - \frac{\alpha}{2}\,F^a_{MN}F^{a\,MN}\Big]~,
\end{equation}
where we have set the AdS radius to one and absorbed the gauge coupling into
the field, $F^a = dA^a + \tfrac12\epsilon^{abc}A^b\wedge A^c$, so that the
Yang--Mills term in its usual normalisation, $-F^2/4\hat g^2$, appears with
the single dimensionless coefficient
\begin{equation}
\alpha = \frac{\kappa^2}{\hat g^2}~.
\end{equation}
This is the only parameter of the theory.\footnote{The backreacted $p$-wave
superfluid of~\cite{Ammon:2009xh} uses the same theory with the unsquared
ratio $\kappa/\hat g$ as its parameter, so its couplings are the square roots
of ours.} In the probe limit $\alpha \to 0$
the gauge field does not deform the geometry, and the calculation reduces to
the one of~\cite{Ammon:2011je}. The equations of motion in Ricci form are
\begin{equation}
R_{MN} + 4\,g_{MN} = \alpha\Big(F^a_{MP}F^{a\phantom{N}P}_{N} - \tfrac16\,g_{MN}F^a_{PQ}F^{a\,PQ}\Big)~, \qquad
D_M F^{a\,MN} = 0~,
\end{equation}
with $D_M$ the gauge-covariant derivative. Since $\alpha$ multiplies the
whole Yang--Mills sector, a background magnetic field $B$ enters the geometry
only through the combination $\beta = \alpha B^2$, and it is this
combination that organises the calculation. Although $\alpha$ is a ratio of
bulk normalisations, it is a boundary observable, proportional to the ratio
$C_J/C_T$ of the current and stress-tensor central charges
(section~\ref{sec:4.4}).

\subsection{The background is exact, and only the condensate is expanded}
\label{sec:2.2}

Near the onset the condensate is small, so a Landau expansion in it is
the right tool as long as the transition is second order; the expansion
itself shows where it is not (section~\ref{sec:5.3}). The backreaction of
the field is not small: $\beta = \alpha B_c^2$ reaches one already at $\alpha \approx 0.03$.
But it need not be treated perturbatively, because the condensate-free
background at any field and temperature can be constructed
non-perturbatively. It is the
magnetic black brane of D'Hoker and Kraus~\cite{DHoker:2009mmn,DHoker:2009ixq},
\begin{equation}
ds^2 = -U\,dt^2 + \frac{dr^2}{U} + e^{2V}\big(dx^2 + dy^2\big) + e^{2Z}dz^2~, \qquad
A^3 = B\,x\,dy~,
\end{equation}
with the flux in the $(x,y)$ plane, which is warped by $V$, and the field
axis along $z$, which is warped separately by $Z$.\footnote{D'Hoker and Kraus write $W$ for our $Z$; we reserve $W$ for the charged gauge field.}
The background field lies in the abelian subgroup generated by $\tau^3$,
so without a condensate the $SU(2)$ theory reduces to Einstein--Maxwell
theory and the brane is exactly the solution of~\cite{DHoker:2009mmn}. The
non-abelian structure enters only through the charged components
$W = A^1 + iA^2$, which see $A^3$ as a $U(1)$ field of unit charge. The
three metric functions obey reduced Einstein--Yang--Mills equations in
which $\alpha$ and $B$ appear only as $\beta$
(appendix~\ref{app:numerics-brane}), so the backgrounds form a
one-parameter family.

We therefore treat the onset, a linear problem on this brane, exactly in
$\beta$, and the lattice perturbatively in the condensate amplitude alone.
Every quantity is then exact in the coupling, at whatever order in the
condensate it is computed. We construct the brane numerically
(appendix~\ref{app:numerics-brane}); appendix~\ref{app:dictionary} gives
the dictionary to the conventions of~\cite{DHoker:2009mmn}.

We work in a fixed-temperature frame throughout: the horizon is at radius
$r_p = 1$ with $U'(1) = 4$, so that $T = 1/\pi$ exactly, and the boundary
normalisation $U, e^{2V}, e^{2Z} \to r^2$ as $r \to \infty$ fixes the
remaining freedom. We quote fields as the dimensionless
$B u_H^2 = B/(\pi T)^2$, where $u_H = 1/\pi T$ is the horizon radius in the
coordinate $u = 1/r$ of the probe literature, so that the probe value
$B_c u_H^2 = 5.13$, that is $B_c \simeq 50\,T^2$, refines the value $5.15$
quoted in~\cite{Ammon:2011je}. At $\beta = 0$ the brane is
AdS$_5$--Schwarzschild, and as $\beta \to \infty$ its interior develops the
long throat of section~\ref{sec:4}.

Table~\ref{tab:notation} collects the notation used in the rest of the
paper.

\begin{table}[t]
\centering
\small
\begin{tabular}{l>{\raggedright\arraybackslash}p{0.62\textwidth}l}
\toprule
symbol & meaning & defined in \\
\midrule
$\alpha$, $\beta$ & coupling $\kappa^2/\hat g^2$; backreaction parameter $\alpha B^2$ & \S\ref{sec:2.1} \\
$U$, $V$, $Z$ & metric functions of the brane; $e^{2V}$ is the proper size of the flux plane & \S\ref{sec:2.2} \\
$B_c u_H^2$ & critical field in horizon units, $B_c/(\pi T)^2$ & \S\ref{sec:2.2} \\
$W$, $w$, $w_0$ & charged field $A^1 + iA^2$; its radial profile; the zero mode at onset & \S\ref{sec:3.1} \\
$P$, $Q$ & Sturm--Liouville coefficients $Ue^Z$, $Be^{Z-2V}$ (written $p_2$, $Bp_1$ in \S\ref{sec:5}) & \S\ref{sec:3.1} \\
$\alpha_\star$, $\alpha_c(3)$ & critical coupling $2/3$; its $d = 3$ counterpart $3.037$ & \S\S\ref{sec:4.2}, \ref{sec:4.3} \\
$\ell_3$, $v$, $b$, $b_h$ & throat radius; frozen $e^{2V}$; throat field $B\ell_3^2e^{-2V}$; $b$ at the horizon & \S\S\ref{sec:4.1}, \ref{sec:4.2} \\
$\varepsilon$ & distance of the horizon invariant from the fixed point, $6 - \beta e^{-4V}$ & \S\ref{sec:4.2} \\
$\rho$, $h$, $\rho_\star$ & radius $r - r_0$ in the throat; BTZ horizon; doubling radius of $e^{2V}$ & \S\S\ref{sec:4.1}, \ref{sec:4.2}, \ref{sec:4.6} \\
$c_0$, $c_1$ & coefficients of the threshold mode, $\rho^{(d-2)/2}w = c_0 + c_1\ln(\rho/\rho_\star)$ & \S\ref{sec:4.2} \\
$\nu$, $\chi$, $\phi_\star$ & $\sqrt{b_h - 1}$; horizon and outer phases & \S\ref{sec:4.6} \\
$C_J$, $C_T$ & current and stress-tensor central charges & \S\ref{sec:4.4} \\
$\eta$, $\epsilon$ & condensate amplitude; coefficient of $\eta^2$, $\propto B - B_c$ & \S\ref{sec:5.1} \\
$G$, $s$, $\gamma$ & reciprocal lattice vector; squared harmonic $G^2$; $s/B$ & \S\ref{sec:5.1} \\
$K$, $K_0$, $K_{G=0}$ & quartic kernel; its limit $s \to 0$; the uniform response & \S\ref{sec:5.1} \\
$C_4$, $X$, $\mathcal{P}$ & contact term; photon screening; metric pairing & \S\ref{sec:5.1} \\
$\tau$, $C(\tau)$ & lattice shape; shell sum selecting it ($C_\triangle$, $C_\square$ at $e^{i\pi/3}$, $i$) & \S\S\ref{sec:5.1}, \ref{sec:5.2} \\
$\alpha_L$, $\alpha_I$ & uniform lattice unstable; $C$ unbounded towards the stripe edge & \S\S\ref{sec:5.3}, \ref{sec:5.2} \\
\bottomrule
\end{tabular}
\caption{Notation. Symbols used only in the appendices are defined there.}
\label{tab:notation}
\end{table}
